\\chapter*{Conclusion générale}
\\addcontentsline{toc}{chapter}{Conclusion générale}
\\markboth{\\textbf{Conclusion générale}}{}

Ce projet de fin d'études nous a permis de mettre en pratique nos connaissances et compétences acquises tout au long de notre formation d'ingénieur en Génie Logiciel et Systèmes d'Information. Grâce à une analyse approfondie des besoins et des contraintes de la société ASTEELFLASH, nous avons pu concevoir et développer une application web complète et fonctionnelle répondant pleinement à leurs attentes en matière de gestion du parc informatique.

\\section*{Bilan du travail réalisé}

L'objectif principal de ce projet était de remplacer le système existant, basé sur des fichiers Excel et des décharges papier, par une application web centralisée et automatisée. Cet objectif a été atteint avec succès grâce à une démarche méthodique articulée autour de la méthodologie agile Scrum et d'un développement itératif en quatre sprints.

L'application développée offre une gamme complète de fonctionnalités essentielles pour la gestion du parc informatique d'ASTEELFLASH :

\\begin{itemize}[label=\\textbullet,font=\\normalsize]
    \\item \\textbf{Authentification sécurisée via LDAP :} L'intégration avec Active Directory garantit une authentification centralisée et une gestion fine des autorisations basée sur les groupes d'utilisateurs.
    \\item \\textbf{Gestion des équipements :} Un inventaire centralisé permet de recenser, suivre et gérer l'ensemble des équipements informatiques (desktops, laptops, serveurs, switchs, imprimantes) avec leurs caractéristiques détaillées.
    \\item \\textbf{Tableau de bord statistique :} Un tableau de bord interactif fournit une vue synthétique de l'état du parc informatique, incluant la répartition par type, par état et par disponibilité.
    \\item \\textbf{Gestion des affectations et missions :} Le système permet d'affecter les équipements aux employés et de gérer les missions temporaires avec un système de rappels automatiques par e-mail.
    \\item \\textbf{Alertes de stock minimum :} Un mécanisme d'alerte automatique notifie les administrateurs lorsque le stock d'un type d'équipement descend en dessous du seuil critique.
    \\item \\textbf{Traçabilité complète :} L'historique détaillé de toutes les opérations (équipements, missions, fournisseurs, commandes) assure une traçabilité totale pour l'audit et le suivi.
    \\item \\textbf{Gestion des commandes et fournisseurs :} Le module de gestion des bons d'achat et des fournisseurs simplifie le processus d'approvisionnement en équipements.
\\end{itemize}

\\section*{Compétences acquises}

Ce projet nous a permis d'acquérir et de renforcer de nombreuses compétences techniques et méthodologiques :

\\begin{itemize}[label=\\textbullet,font=\\normalsize]
    \\item Maîtrise du framework ASP.NET Core et du patron architectural MVC pour le développement d'applications web professionnelles.
    \\item Utilisation d'Entity Framework Core comme ORM pour la gestion de la persistance des données et l'abstraction de la couche d'accès aux données.
    \\item Intégration avec Microsoft Active Directory via le protocole LDAP pour l'authentification et la gestion des autorisations.
    \\item Configuration et utilisation de Microsoft SQL Server 2019 pour la conception et la gestion de bases de données relationnelles.
    \\item Application pratique de la méthodologie agile Scrum dans un contexte professionnel réel, avec des sprints, des revues et une collaboration continue avec le Product Owner.
    \\item Modélisation UML (diagrammes de cas d'utilisation, de classes, de séquences et de déploiement) pour la conception logicielle.
    \\item Mise en place de services de messagerie SMTP pour l'envoi automatique de notifications et d'alertes.
\\end{itemize}

\\section*{Limites et perspectives}

Bien que l'application réponde aux objectifs fixés, certaines améliorations peuvent être envisagées pour les versions futures :

\\begin{itemize}[label=\\textbullet,font=\\normalsize]
    \\item \\textbf{Gestion des licences logicielles :} Ajouter un module dédié au suivi des licences logicielles, incluant les dates d'expiration, les renouvellements et les coûts associés.
    \\item \\textbf{Inventaire automatisé :} Intégrer un agent d'inventaire automatique (type OCS Inventory) pour détecter et recenser automatiquement les équipements connectés au réseau de l'entreprise.
    \\item \\textbf{Système de ticketing :} Développer un module de gestion des incidents et des demandes de support, permettant aux utilisateurs de signaler les problèmes techniques et de suivre leur résolution.
    \\item \\textbf{Monitoring réseau :} Intégrer des fonctionnalités de surveillance en temps réel de l'état des équipements réseau (serveurs, switchs, points d'accès) avec des alertes proactives.
    \\item \\textbf{Application mobile :} Développer une version mobile de l'application pour permettre aux techniciens IT de gérer le parc informatique lors de leurs déplacements dans l'usine.
    \\item \\textbf{Reporting avancé :} Enrichir les fonctionnalités de reporting avec des rapports personnalisables, des exports PDF et des graphiques analytiques pour faciliter la prise de décisions stratégiques.
    \\item \\textbf{Gestion financière des actifs :} Ajouter un module de suivi financier incluant l'amortissement des équipements, le coût total de possession (TCO) et la planification budgétaire.
\\end{itemize}

Notre application est conçue pour être évolutive et adaptable aux mises à jour futures. L'architecture modulaire basée sur ASP.NET Core et le patron MVC facilite l'ajout de nouvelles fonctionnalités sans impact sur les modules existants. Nous sommes convaincus que cette solution constitue une base solide pour accompagner la croissance et l'évolution des besoins informatiques d'ASTEELFLASH.

En définitive, ce projet de fin d'études a été une expérience enrichissante tant sur le plan technique que professionnel. Il nous a permis de contribuer concrètement à l'amélioration de la gestion du patrimoine informatique d'ASTEELFLASH tout en consolidant nos compétences en développement web et en gestion de projet.
